\chapter{개요}

\section{과제 목적과 연구 질문}

800 V급 배터리는 같은 전력을 더 작은 전류로 전달해 급속충전과 배선 손실에 유리하다\cite{b1}. 그러나 차량 안팎의 전압은 하나로 통일되지 않는다. 차량은 다음 기능을 함께 수행해야 한다.

\begin{itemize}
\item 400·500 V급 충전설비와의 호환
\item 400 V급 보조부하(컴프레서, 히터)에 대한 전원 공급
\item 48 V·12 V 전원망의 절연 변환과 고장 시 공급 연속성
\item 여러 부하가 겹칠 때 변환기 정격을 넘지 않도록 하는 전력 관리
\end{itemize}

전전력 절연 변환기는 정밀 제어가 쉽지만 150 kW 전부를 처리하는 자기소자, 반도체와 냉각이 필요하다\cite{b2}. 2:1 고정비 변환기는 효율이 높지만\cite{b3} 출력 전압이 배터리 전압을 따라 움직인다. PRISM-Hub는 대전력 전달과 조정을 분리해 이 상충을 다룬다. 고정비 경로와 소용량 부분전력 조정단의 결합, 직렬 전압보상은 선행연구에 원리가 있다\cite{b4,b5,b6}. 본 과제는 이 원리를 차량의 충전, 400 V 보조버스와 48 V·12 V 전원망에 배치하고, 포트 정격·운전 모드·고장 강등을 하나의 시스템으로 검증하는 데 범위를 둔다. 연구 질문은 다음 세 가지다.

\begin{itemize}
\item Q1(정격 분리). 150 kW급 충전전력을 조정단에서 분리한 상태에서, 소용량 조정단이 보조버스의 차동전력만으로 배터리 전 범위를 조정할 수 있는가.
\item Q2(강건 제어와 부하 관리). 하나의 제어기와 인과적 부하 승인 계층이 부품 공차, 디지털 제어 지연, 부하 중첩과 고장 조건에서 포트 정격과 전압 품질을 함께 지킬 수 있는가.
\item Q3(기준 구조 대비 이점). 이 구조가 조정 없는 고정비 버스나 전전력 변환기를 쓰는 구조보다 손실, 정격 또는 부품 수에서 측정 가능한 이점을 주는가. 준다면 어떤 팩 등급과 부품 등급에서인가.
\end{itemize}

세 질문은 다섯 가설(H1–H5)과 수락 기준으로 판정한다(3.2절). Q3의 답이 사양 팩에서 부정적으로 나와(6장) 구조를 개선했으므로, Q3에는 적용 범위(7장)와 산업에서 쓰는 방식 대비 성립 조건(15장)을 묻는 부분을 더했다. 이번 과제의 범위는 MATLAB/Simulink와 Python을 이용한 시뮬레이션 기반 설계 검증과 데이터시트 대조이며, 하드웨어 실증은 후속 단계에서 수행한다.

\section{보고 범위와 중간보고서 대비 변경 사항}

본 보고서는 PRISM-Hub의 설계·검증 결과 전체를 다룬다. 중간보고서의 수치 가운데 Stage A 손실·충전·클램프 값은 150 kHz 원설계의 값이라, 이 보고서에서는 최종 설계(A25R) 값을 쓴다. 바뀐 값과 그 이유는 14장에 모았다. 중간보고서 대비 주요 변경은 표 \ref{t:changes}과 같다.

{\footnotesize
\begin{longtable}{L{\dimexpr 0.18\linewidth-2\tabcolsep\relax}L{\dimexpr 0.28\linewidth-2\tabcolsep\relax}L{\dimexpr 0.44\linewidth-2\tabcolsep\relax}L{\dimexpr 0.10\linewidth-2\tabcolsep\relax}}
\caption{중간보고서(2026-10-26) 대비 주요 변경 사항}\label{t:changes}\\
\toprule
\tabhead \textbf{항목} & \textbf{중간보고서} & \textbf{최종 보고서} & \textbf{관련 장} \\
\midrule
\endfirsthead
\toprule
\tabhead \textbf{항목} & \textbf{중간보고서} & \textbf{최종 보고서} & \textbf{관련 장} \\
\midrule
\endhead
\bottomrule
\endlastfoot
Stage A 설계 & 150 kHz, 위치당 6 mΩ 가정, V\textsubscript{A} 노드 커패시터 400~µF & A25R: 25 kHz, 위치당 C4MS008075U2 2병렬, L\textsubscript{r} 864.6 nH, C\textsubscript{fly} 46.875 µF(부품 45.0 µF), C\textsubscript{A} 2.4 mF. 출력용량 손실 모델 정정(E31–E35, TB25–TB27) & 13 \\
손실 비교와 창 조정 모듈 & ④가 ⑤보다 26 \% 작음. B\textsubscript{420} 200–216s, B\textsubscript{470} 224–240s에서 창 조정 모듈이 최소 손실 & ④와 ⑤는 보조 손실이 같은 수준(공정 기준 −1.4 \%), ④는 ⑥보다 4–33 W 큼. 공정 기준에서 창 조정 모듈은 같은 칸의 다른 성립 구조와 보조 손실이 같은 수준이라 성립 수단으로만 위치(표 \ref{t:recalc}; E41–E45) & 6, 7, 14, 15 \\
보호·부하 관리 & 보호 상태기계 단독 & SALS 분리·감발 연동, 48 V 회생 차단, Stage C 보호 일시 정지, 이중 고장 400 V 부하 정지 3 ms(E29, E38) & 9 \\
충전 시작·정지와 클램프 & 충전전류 증가 램프 100~µs 이상. 클램프 1 Ω·7 kJ, 회로 3/4(방향 반전 1건 불통과) & 사전 기동과 램프 600 µs 이상(E35; 보조부하 15·25 kW에서 확인, E46). 방향 반전 처리(E27, TB21) 뒤 부호 없는 영교차 게이트(E28). 클램프 가지 3개 × 2.4 Ω, 최대 연속 켜짐 0.5 ms, 구현 요구는 표 \ref{t:e26r}(E38–E40, E44, E45) & 8, 11, 12, 13 \\
선행기술·기여 & 12건 비교, "같은 2:1 단의 계속 사용은 찾지 못함" & 29건 대조, 개념의 선행이 있어 기여 1을 "성립 조건과 같은 노드 배치의 회로 검증"으로 고침 & 15 \\
\end{longtable}
}
